\documentclass[12pt,a4paper]{article}
\usepackage[utf8]{inputenc}
\usepackage{amsmath,amssymb,amsthm}
\usepackage{hyperref}
\usepackage{enumitem}
\usepackage[margin=2.5cm]{geometry}

\newtheorem{theorem}{Theorem}[section]
\newtheorem{lemma}[theorem]{Lemma}
\newtheorem{definition}[theorem]{Definition}
\newtheorem{corollary}[theorem]{Corollary}

\title{BCRD Operator Algebra:\\[0.3em]
Trace-Preserving Involutive Evolution on Six-Index Matrices\\[0.5em]
\large A $*$-Automorphism from the Frozen Unitary}
\author{Tobias Croydon-McRae\\
\textit{Independent Researcher, Aotearoa New Zealand}\\
\texttt{github.com/Tobillicious/universal-protection-theory}}
\date{November 2026}

\begin{document}
\maketitle

\begin{abstract}
We define the BCRD evolution operator $\Phi(A) = W A W^\top$ on
$6 \times 6$ integer matrices, where $W$ is a frozen permutation
unitary (the M-sector $2 \leftrightarrow 3$ swap). We prove that
$\Phi$ constitutes a $*$-automorphism of the matrix algebra:
it preserves trace, matrix product, transpose (adjoint), and the
identity matrix, and is involutive ($\Phi^2 = \mathrm{id}$). The
master lemma --- that conjugation by $W$ merely permutes row and
column indices --- is proved by exhaustive finite case split over
$\mathrm{Fin}\, 6 \times \mathrm{Fin}\, 6$, reducing every
algebraic law to a reindexing. We additionally prove the
Hilbert--Schmidt inner product is preserved and establish the
diagonal correspondence connecting the operator to the independently
specified BCRD source state generator. All theorems verified in
Lean~4 with standard axioms only. Zero \texttt{sorry}.
\end{abstract}

\section{Introduction}

The BCRD programme models the dynamics of a six-position relational
system. The evolution operator acts on $6 \times 6$ matrices by
conjugation with a frozen unitary~$W$, encoding the M-sector swap
(positions $2 \leftrightarrow 3$). This paper proves that this
conjugation is a full $*$-automorphism of the matrix algebra ---
the strongest algebraic characterization available.

The operator is \emph{not} an arbitrary similarity transform. It is
a permutation matrix (hence orthogonal, hence unitary over the reals),
and its algebraic properties follow entirely from the permutation
structure. No ring lemmas, norm\_num, or mathlib are required; the
proofs proceed by finite decidable case splits.

\section{Definitions}

\begin{definition}[Matrix6]
A $6 \times 6$ integer matrix: $\text{Matrix6} := \mathrm{Fin}\, 6 \to
\mathrm{Fin}\, 6 \to \mathbb{Z}$.
\end{definition}

\begin{definition}[The Frozen Permutation]
The permutation $\sigma : \mathrm{Fin}\, 6 \to \mathrm{Fin}\, 6$
swaps positions $2$ and $3$ and fixes the rest:
\[
\sigma(0) = 0, \quad \sigma(1) = 1, \quad \sigma(2) = 3, \quad
\sigma(3) = 2, \quad \sigma(4) = 4, \quad \sigma(5) = 5
\]
\end{definition}

\begin{definition}[The Unitary $W$]
$W_{ij} = 1$ if $j = \sigma(i)$, else $0$. This is the permutation
matrix of~$\sigma$.
\end{definition}

\begin{definition}[Evolution Operator]
$\Phi(A) := W \cdot A \cdot W^\top$ (conjugation by $W$).
\end{definition}

\section{The Master Lemma}

\begin{theorem}[Conjugation Permutes Indices]\label{thm:master}
For all $A : \text{Matrix6}$ and indices $i, j : \mathrm{Fin}\, 6$:
\[
\Phi(A)_{ij} = A_{\sigma(i),\, \sigma(j)}
\]
\end{theorem}

\begin{proof}
By exhaustive case split over $\mathrm{Fin}\, 6 \times \mathrm{Fin}\, 6$
(36 cases). Each case reduces to arithmetic on the $W$ entries via
\texttt{wUpdate\_apply} (the lever lemma: $W_{ij} = [j = \sigma(i)]$)
and \texttt{simp} with \texttt{decide}.
\end{proof}

\begin{theorem}[Permutation Is Involutive]\label{thm:invol-perm}
$\sigma \circ \sigma = \mathrm{id}$. That is, $\forall\, x,\;
\sigma(\sigma(x)) = x$. By \texttt{decide}.
\end{theorem}

\section{The $*$-Automorphism}

All results below follow from the master lemma by reindexing.

\begin{theorem}[Trace Preservation]\label{thm:trace}
$\mathrm{tr}(\Phi(A)) = \mathrm{tr}(A)$.
\end{theorem}

\begin{proof}
$\mathrm{tr}(\Phi(A)) = \sum_i \Phi(A)_{ii} = \sum_i A_{\sigma(i),\sigma(i)}
= \sum_i A_{ii}$ (the permutation is a bijection on the diagonal).
\end{proof}

\begin{theorem}[Product Preservation]\label{thm:product}
$\Phi(A \cdot B) = \Phi(A) \cdot \Phi(B)$ (algebra homomorphism).
\end{theorem}

\begin{proof}
$[\Phi(AB)]_{ij} = (AB)_{\sigma(i),\sigma(j)} = \sum_k A_{\sigma(i),k}
B_{k,\sigma(j)}$. Reindex $k = \sigma(m)$:
$= \sum_m A_{\sigma(i),\sigma(m)} B_{\sigma(m),\sigma(j)}
= \sum_m \Phi(A)_{im} \Phi(B)_{mj} = [\Phi(A)\Phi(B)]_{ij}$.
\end{proof}

\begin{theorem}[Transpose Preservation]\label{thm:transpose}
$\Phi(A^\top) = (\Phi(A))^\top$ (adjoint preservation).
\end{theorem}

\begin{proof}
$[\Phi(A^\top)]_{ij} = A^\top_{\sigma(i),\sigma(j)} = A_{\sigma(j),\sigma(i)}
= \Phi(A)_{ji} = [(\Phi(A))^\top]_{ij}$.
\end{proof}

\begin{theorem}[Identity Preservation]\label{thm:identity}
$\Phi(I) = I$.
\end{theorem}

\begin{proof}
$\Phi(I)_{ij} = I_{\sigma(i),\sigma(j)} = [\sigma(i) = \sigma(j)]
= [i = j] = I_{ij}$ (the permutation is injective).
By \texttt{decide}.
\end{proof}

\begin{theorem}[Involutivity]\label{thm:involutive}
$\Phi(\Phi(A)) = A$ ($\Phi$ is its own inverse).
\end{theorem}

\begin{proof}
$[\Phi(\Phi(A))]_{ij} = \Phi(A)_{\sigma(i),\sigma(j)}
= A_{\sigma(\sigma(i)),\sigma(\sigma(j))} = A_{ij}$
by Theorem~\ref{thm:invol-perm}.
\end{proof}

\section{Inner Products}

\begin{theorem}[Trace Bilinear Form]\label{thm:trace-form}
$\mathrm{tr}(\Phi(A) \Phi(B)) = \mathrm{tr}(A B)$.
\end{theorem}

\begin{proof}
$\mathrm{tr}(\Phi(A)\Phi(B)) = \mathrm{tr}(\Phi(AB))
= \mathrm{tr}(AB)$ by Theorems~\ref{thm:product} and~\ref{thm:trace}.
\end{proof}

\begin{theorem}[Hilbert--Schmidt / Frobenius Inner Product]\label{thm:hs}
$\langle \Phi(A), \Phi(B) \rangle_{\mathrm{HS}} :=
\mathrm{tr}((\Phi(A))^\top \Phi(B)) = \mathrm{tr}(A^\top B)$.
$\Phi$ is an isometry for the Hilbert--Schmidt metric.
\end{theorem}

\begin{proof}
By Theorems~\ref{thm:transpose}, \ref{thm:product}, and~\ref{thm:trace}.
\end{proof}

\section{Diagonal Correspondence}

\begin{theorem}[Evolve Diagonal]\label{thm:diag}
For every BCRD source state configuration $s : \mathrm{Fin}\, 6 \to \mathbb{Z}$:
\[
\Phi(\mathrm{diag}(s)) = \mathrm{diag}(\text{step}(s))
\]
where $\text{step}$ is the independently specified BCRD generator action
(the M-sector transport from \texttt{BCRDActualRelationalFrame}).
\end{theorem}

This connects the arbitrary-configuration matrix operator to the
independently specified physical generator, generalizing the single-event
demonstration \texttt{evolve\_realises\_event}.

\section{Axiom Audit}

\begin{itemize}
    \item \texttt{wUpdate\_apply}: no axioms (by \texttt{decide})
    \item \texttt{permutation\_involutive}: no axioms (by \texttt{decide})
    \item \texttt{evolve\_entry}: no axioms
    \item \texttt{evolve\_preserves\_trace}: no axioms
    \item \texttt{evolve\_preserves\_product}: no axioms
    \item \texttt{evolve\_preserves\_transpose}: no axioms
    \item \texttt{evolve\_preserves\_identity}: no axioms
    \item \texttt{evolve\_involutive}: no axioms
    \item \texttt{evolve\_preserves\_frobenius}: no axioms
    \item \texttt{evolve\_diagonal}: \texttt{propext}, \texttt{Quot.sound}
\end{itemize}

\section{Verification}

Verified in \texttt{BCRDOperatorAlgebra.lean} with imports from
\texttt{FrozenBCRDEvents.lean} and \texttt{BCRDActualRelationalFrame.lean}.
Build: 496 jobs, zero errors, zero \texttt{sorry}.

\section{Conclusion}

The BCRD evolution operator is a full $*$-automorphism of the
$6 \times 6$ integer matrix algebra: trace-preserving, product-preserving,
adjoint-preserving, identity-preserving, and involutive. It is an
isometry for both the trace form and the Hilbert--Schmidt inner product.
The diagonal correspondence connects it to the independently specified
BCRD source state generator. This is the algebraic backbone of the
BCRD programme.

\end{document}
